\chapter{File formats}
\label{chap:formats}

AirfoilTools can read and write several airfoil formats, which
fall into two families: the \emph{discrete points} formats
(originating from the aeronautics field) and the \emph{spline}
formats specific to AirfoilTools.

\section{Discrete points formats}

\subsection{Selig format (\chemin{.dat})}

The most widespread historical format, popularized by the
University of Illinois (UIUC) database. It is the default format
on \chemin{airfoiltools.com}.

\begin{itemize}
    \item First line: airfoil name (free).
    \item Following lines: \texttt{x y} pairs separated by
        spaces.
    \item Traversal in \emph{Selig convention}: from the trailing edge
        ($x=1$) towards the upper surface (decreasing x), to the leading
        edge ($x=0$), then towards the lower surface (increasing x)
        and back to the trailing edge.
\end{itemize}

\begin{lstlisting}[caption={Example Selig file (NACA~2412)}]
NACA 2412
   1.000000   0.001260
   0.985700   0.003930
   0.971440   0.006580
   ...
   0.014300   0.018780
   0.000000   0.000000
   0.014300  -0.014050
   ...
   0.985700  -0.001280
   1.000000  -0.001260
\end{lstlisting}

\subsection{Lednicer format (\chemin{.dat})}

A variant of Selig where upper surface and lower surface are stored
separately, each traversed from the leading edge towards the trailing
edge. A header line indicates the number of points on each
side.

\begin{lstlisting}[caption={Example Lednicer file}]
NACA 2412
65.0   65.0

   0.000000   0.000000
   0.001425   0.005700
   ...
   1.000000   0.001260

   0.000000   0.000000
   0.001425  -0.005720
   ...
   1.000000  -0.001260
\end{lstlisting}

\subsection{CSV format (\chemin{.csv})}

A tabular format with header \chemin{x,y}. Convenient for
importing into spreadsheets or for sharing with
non-aeronautical tools.

\begin{lstlisting}[caption={Example CSV file}]
x,y
1.000000,0.001260
0.985700,0.003930
...
\end{lstlisting}

\subsection{GNU format (\chemin{.gnu})}

A raw points format compatible with the \textbf{Axile} software and
directly usable by \textbf{gnuplot}. It is a simple
text file of two columns \texttt{x~y}, without header or airfoil name.

\begin{itemize}
    \item One line per point, two columns \texttt{x y}.
    \item Each number in scientific notation with 18 decimals
        (\texttt{\%.18e}), separated by a single space.
    \item \texttt{LF} line endings (including after the last point),
        guaranteed even on Windows.
    \item No header, comment or third column.
\end{itemize}

\begin{lstlisting}[caption={Example GNU file (.gnu)}]
1.000000000000000000e+00 0.000000000000000000e+00
9.754952615109063752e-01 6.041608209341215764e-03
0.000000000000000000e+00 0.000000000000000000e+00
\end{lstlisting}

This format is produced identically by \code{numpy.savetxt} and
readable back by \code{numpy.loadtxt} (array of shape $(n, 2)$).

\begin{noteinfo}
The \chemin{.gnu} format is offered as \textbf{write-only}
(save of the current airfoil). To re-import points,
use a \chemin{.dat} (Selig) or \chemin{.csv} file.
\end{noteinfo}

\section{AirfoilTools spline formats}

\subsection{\chemin{.bspl} format}

A complete save of an airfoil in \textbf{multi-segment
spline} mode. It preserves the entire Bézier definition:
control points, degrees, continuities at junctions, sampling
parameters.

This is the \textbf{recommended} format to keep the editing
work on an airfoil and re-open it without loss of information.

\begin{noteinfo}
The \chemin{.bspl} format is a structured text format (JSON
or similar) specific to AirfoilTools. It is not
interoperable with other CAD or aerodynamics software.
For external distribution, re-export to
\chemin{.dat} (Selig) after adjusting the sampling.
\end{noteinfo}

\subsection{\chemin{.bez} format}

A simplified variant of the \chemin{.bspl} for airfoils defined
by a single Bézier curve per side (single-segment mode).
Kept for compatibility with earlier versions.

\section{Project format (\chemin{.aftproj})}
\label{sec:aftproj}

The \chemin{.aftproj} format saves the \textbf{entire working
session} in a single file, where the preceding formats
only save an isolated airfoil. A project contains:

\begin{itemize}
    \item the \textbf{current airfoil} and the \textbf{reference
        airfoil} (with their complete spline definition where
        applicable);
    \item the optional \textbf{tracing image}, \emph{embedded}
        in the file, together with its \textbf{transformation}
        (position, scale, rotation, visibility).
\end{itemize}

It is a \textbf{JSON} file (text). The image is encoded
in base64 and embedded directly: the project is therefore
\textbf{self-contained and portable}. The tracing image is recovered
on reopening even if the original image file has been moved or
deleted.

\begin{itemize}
    \item Open: \menu{File $\rightarrow$ Open project\dots}
        (\touche{Ctrl}+\touche{Maj}+\touche{P}).
    \item Save: \menu{File $\rightarrow$ Save
        project\dots} (\touche{Ctrl}+\touche{Alt}+\touche{S}).
\end{itemize}

\begin{important}
The \chemin{.aftproj} format is the \textbf{only} one that keeps
the tracing image and its alignment. If your work involves
tracing (see section~\ref{sec:calque}), save a project
rather than a simple airfoil.
\end{important}

\begin{noteinfo}
Since the image is embedded in base64, a project containing a large
image produces an \chemin{.aftproj} file larger than the
sum of the airfoil files. This is the price of portability
(no dependency on an external image file).
\end{noteinfo}

\section{NACA generation}

The application can generate NACA airfoils on the fly from
their designation, without going through a file:

\begin{itemize}
    \item NACA 4 digits (e.g. \texttt{2412}, \texttt{0012}):
        maximum camber (\%), position of the max camber
        (tenths), relative thickness (\%).
    \item NACA 5 digits (e.g. \texttt{23012}): cambered airfoil
        with reflex mean line.
\end{itemize}

This generation happens automatically at application
startup (NACA~2412 as the current airfoil and as the reference).
To generate other NACA airfoils during a session, use the
\menu{File $\rightarrow$ NACA airfoil $\rightarrow$ Current
airfoil\dots} or \menu{\dots Reference airfoil\dots} menu: a dialog
asks for the indices, then the number of points. See the detailed
procedure in section~\ref{sec:naca}, page~\pageref{sec:naca}.

\begin{noteinfo}
A NACA airfoil not yet modified can be cleanly regenerated at
another resolution via the \menu{Number of
points\dots} context menu (right-click), which recalls the analytical
formula instead of interpolating the points. See section~\ref{sec:naca}.
\end{noteinfo}

\section{Summary table}

\begin{tabularx}{\linewidth}{l c c X}
    \toprule
    \textbf{Format} & \textbf{Read} & \textbf{Write}
        & \textbf{Preserves the spline?} \\
    \midrule
    Selig (\chemin{.dat})    & yes & yes & no \\
    Lednicer (\chemin{.dat}) & yes & yes & no \\
    CSV (\chemin{.csv})      & yes & yes & no \\
    GNU (\chemin{.gnu})      & no & yes & no \\
    Spline (\chemin{.bspl})  & yes & yes & \textbf{yes} \\
    Bézier (\chemin{.bez})   & yes & yes & \textbf{yes} (single-segment) \\
    Project (\chemin{.aftproj}) & yes & yes & \textbf{yes} (+ tracing) \\
    \bottomrule
\end{tabularx}

\begin{important}
When you save an airfoil in \chemin{.dat} or
\chemin{.csv} format, \textbf{the spline information is lost}:
only the list of sampled discrete points is written.
To fully preserve your work, always save
in \chemin{.bspl}.
\end{important}
